
\section{Introduction}

This article studies the metric theory of Diophantine approximation along a
\emph{sparse, super-exponentially growing} subsequence of denominators: the
selected convergent denominators $(Q_j)_{j\ge1}$ of $\ee$ (notation in the
Setup below). We first show in Theorem~\ref{thm:full-spectrum} that the integer
three-term recurrence for $(Q_j)$ forces complete discretization of the weighted
approximation spectrum:
\[
\Lam(\R)=\frac14\Z_{\ge0}\cup\{\infty\}
=\left\{0,\frac14,\frac12,\frac34,1,\dots\right\}\cup\{\infty\},
\]
creating an infinite family of equidistant empty spectral gaps
$(k/4,(k+1)/4)$ of width $1/4$ and forcing the sublevel sets
$E_c=\{x:\Lam(x)\le c\}$ to be constant across each gap
($E_c=E_{k/4}$ for $c\in[k/4,(k+1)/4)$, Corollary~\ref{thm:sub}).

Our main metric result at the critical weight, Theorem~\ref{thm:gauge},
measures each positive quantized level of $\Lam$: for each integer $k\ge1$
there is an explicit Hausdorff gauge $h_k$ with $h_k(1/Q_n)=(2k+1)^{-n}$, built
from the recurrence itself, for which the exact level $\{\Lam=k/4\}$ has
infinite but $\sigma$-finite Hausdorff measure $\mathcal H^{h_k}$ on the circle
$\R/\Z$, while the strict sublevel $\{\Lam<k/4\}$ is $\mathcal H^{h_k}$-null.
This differs both from the $h$-sets of Cabrelli, Mendivil, Molter and
Shonkwiler \cite[\S5, Theorem~3]{CMMS}, whose gauge gives a finite positive
measure, and from the classical Lagrange level at value $3$ studied by Moreira,
Erazo and Angelini \cite{MEA}, which has either zero or non-$\sigma$-finite
measure for every dimension function. Corollary~\ref{cor:exact-gauge} shows
that $h_k$ is the exact critical gauge of the level $k/4$: smaller gauges give
measure zero even on the non-strict sublevel $\{\Lam\le k/4\}$, larger gauges
destroy $\sigma$-finiteness on $\{\Lam=k/4\}$, and the $\mathcal H^{h_k}$
measure of $\{\Lam=k/4\}$ is infinite in every nonempty open arc.
Corollary~\ref{cor:liminf-gauges} classifies every nondecreasing gauge by
its lower limit along the recurrence scales, even when the normalized
cost oscillates. Theorem~\ref{thm:finite-state-tail} shows how this
classification changes under constraints on consecutive defects.
The largest Perron root is taken only over components in which a path
can revisit an extreme defect infinitely often.
Furthermore, Theorem~\ref{thm:trichotomy} establishes a sharp three-way
splitting at each positive spectral level $k/4$ ($k\ge1$) comparing the limsup
level $\{\Lam=k/4\}$ (gauge base $2k+1$), the unsigned exact limit set
$\{x:j\|Q_jx\|\to k/4\}$ (universal critical gauge base $2$ for \emph{all}
$k\ge1$), and the signed exact limit sets
$\{x:j(Q_jx-\mathrm{round}(Q_jx))\to\pm k/4\}$ (each of which collapses to a
single countable $L$-coset). Theorems~\ref{thm:general-rec}--\ref{thm:poly-phase-transition}
extend the quantization, reference gauges, three-way splitting and
dimension results to general polynomial-growth recurrences
$Q_{j+1}=(cj^r+d)Q_j+Q_{j-1}$ ($r,c\in\Z_{\ge1}$, $d\in\Z$), replacing the
$\ee$-convergent lattice $L$ by an intrinsic rank-$2$ lattice $L^{(r,c,d)}$
spanned by fundamental solutions and establishing the discretized spectrum
$\frac1c\Z_{\ge0}\cup\{\infty\}$, exact gauges with $1/r$-scaled logarithmic
asymptotics, positive-level trichotomy, and subcritical dimension phase
transition $1-\alpha/r$ ($0\le\alpha<r$). In Section~\ref{sec:growing-rec}
(Theorems~\ref{thm:growing-rec-spectrum} and~\ref{thm:growing-rec-dim-phase},
and Corollary~\ref{cor:growing-rec-three-regimes}), we remove the polynomial
growth hypothesis altogether and establish the complete spectral, gauge, and
Hausdorff-dimension phase diagram for arbitrary integer recurrences
$Q_{j+1}=A_jQ_j+Q_{j-1}$ with $A_j\ge2$ and $A_j\to\infty$, uncovering both the
universal subcritical distortion law $\dimH=(1-\tau)/(1+\tau\gamma_{\sup})$ and
a secondary supercritical phase transition at $\tau_{\mathrm{crit}}=1+1/\gamma_{\sup}$
where the level sets collapse from the rank-$2$ lattice $L_G$ to the rank-$1$
rational sublattice $\frac{1}{g}\Z$.

On the coarse scale of Hausdorff dimension, Theorem~\ref{thm:finite} shows that
every finite level set $E_c$ has Hausdorff dimension $0$, so $\{\Lam<\infty\}$
has dimension $0$ while its complement $\{\Lam=\infty\}$ has full Lebesgue
measure and Hausdorff dimension $1$. At subcritical weights
$\Lam_\alpha(x)=\limsup_{j\to\infty}j^\alpha\|Q_jx\|$ ($0\le\alpha<1$),
Theorems~\ref{thm:upper} and~\ref{conj:weight} give exact Hausdorff dimension $1-\alpha$ for
$\{x:\Lam_\alpha(x)\le c\}$ for every $c\ge0$, and for
$\{x\in[0,1]:j^\alpha\|Q_jx\|\to c\}$ for every admissible target $c$.
For $\alpha=0$ such targets are $0\le c\le 1/2$.
The proof uses a covering bound and a shifted-grid Moran construction;
the case $\alpha=0$ also follows from Bugeaud, Liao and Rams \cite{BLR}, and
both bounds for $0\le\alpha<1$ are machine-checked in Lean~4. Finally, in
Section~\ref{sec:windows-residues}, Theorems~\ref{thm:windows}--\ref{thm:no-short-dim}
determine the Hausdorff dimensions and almost-sure frequency of Hermite
continued-fraction window constraints, Theorem~\ref{thm:residues} proves
that eventual periodicity of ordinary continued-fraction digits modulo $4$ has
the same Hausdorff dimension as the fixed alphabet $1\bmod 4$, bounded strictly
between $357/512$ and $179/256$, Theorem~\ref{thm:residues-general} extends this
to every modulus $m\ge2$ (via a uniform $24$-Lipschitz sandwich, an unconditional
transfer-operator Hausdorff-dimension framework, and exact integer certificates
for $m\in\{2,3,4,5,6\}$), and Theorem~\ref{thm:residues-asymptotic} proves the
global strict separation $\dimH F_{m,m}<\dimH F_{m,1}=\dimH E_m$ for all $m\ge2$
together with the exact two-term large-$m$ asymptotic expansion governed by the
reciprocal Fibonacci constant $\psi_{\mathrm{Fib}}=\sum_{k=1}^\infty 1/F_k$ and
its metallic-continuant generalizations $C_r=\sum_{k=1}^\infty 1/q_k(r)$.

These results build on the zero-level classification and initial spectral gap
(Section~\ref{sec:class}), which arose as machinery in a separate formalization
project on the (open) irrationality of $e+\pi$ \cite{Phase1}.

\subsection{Setup}

The continued fraction of $\ee$ is $[2;1,2,1,1,4,1,1,6,\dots]$. Let $P_j/Q_j$
be the convergent immediately preceding the partial quotient $2j$; the first
few are $3/1,19/7,193/71,2721/1001$. The denominators satisfy
\begin{equation}\label{eq:rec}
Q_{j+1}=(4j+2)Q_j+Q_{j-1}\quad(j\ge2),\qquad Q_1=1,\ Q_2=7,
\end{equation}
so $Q_{j+1}/Q_j\sim 4j$ and $\log Q_j\sim j\log j$ (super-exponential growth).
Write $\|y\|=\dist(y,\Z)$ and define
\[
\Lam(x):=\limsup_{j\to\infty} j\,\|Q_j x\|\in[0,\infty].
\]
More generally, for $\alpha\ge0$ put
$\Lam_\alpha(x):=\limsup_{j\to\infty} j^\alpha\|Q_j x\|\in[0,\infty]$,
so that $\Lam=\Lam_1$.
Let $L:=\Z+\Z\frac{\ee+1}{2}$, a countable dense subgroup of $\R$.

\subsection{Relation to prior work}\label{sec:prior}

The \emph{phenomenon} (a spectral gap plus zero-dimensional sub-threshold level
sets) is classical: in the classical Lagrange spectrum, the spectrum below
$\sqrt5$ (and below $3$) is countable (Hausdorff dimension $0$), and Hall's ray
gives a half-line. Moreira \cite{Moreira} studies the Hausdorff dimension of
portions of the classical Markov and Lagrange spectra below a threshold. Lima
and Moreira \cite{LimaMoreira} prove phase transitions for the \emph{dynamical}
Markov and Lagrange spectra of generic horseshoes; the corresponding question
for the classical spectra is left open there. Note that our level sets $E_c$ are
subsets of the space of inputs $x\in\R$.

The closest prior continued-fraction result is a theorem of Kraaikamp and
Liardet \cite{KL}, quoted as Lemma~4.1 in \cite{BDS}. Let $\alpha$ be irrational
with convergents $p_n/q_n$. Then a real $\beta$ lies outside
$\Z+\Z\alpha$ if and only if
$\|q_n\beta\|>\tfrac14\,q_n|q_n\alpha-p_n|$ for infinitely many $n$. This
is a classification together with a gap, with constant $1/4$, for the
normalized quantity $\|q_n\beta\|/(q_n|q_n\alpha-p_n|)$ along \emph{all}
convergent denominators. Theorems~\ref{thm:class} and~\ref{thm:gap} can be
compared with this criterion in two complementary ways:
\begin{itemize}[leftmargin=1.8em]
\item \emph{As a subsequence along the convergents of $\ee$:} Passing from all
convergent denominators of $\ee$ to the subsequence $(Q_j)$ of every third
convergent denominator changes two things. First, the zero group grows from
$\Z+\Z\ee$ to $L$: for $x=(\ee+1)/2$ the quantity $j\|Q_jx\|$ tends to $0$, so
the infinitely many $n$ supplied by Kraaikamp and Liardet for $\alpha=\ee$ lie
outside the subsequence $(Q_j)$. Second, the normalization changes. Since
$2j+1<1/(Q_j|Q_j\ee-P_j|)<2j+2$, one has
$\limsup_{j\to\infty}\|Q_jx\|/(Q_j|Q_j\ee-P_j|)=2\Lam(x)$. Thus the gap
$\Lam\ge1/4$ reads $\limsup\ge1/2$ in their normalization, and the value
$\Lam(x_0)=1/4$ is attained by the explicit point $x_0$ in
Remark~\ref{rem:unweighted}. The two constants $1/4$ therefore refer to
different normalizations.
\item \emph{As the full convergent sequence of $\beta=(\ee+1)/2$:} Put
$R_j=(P_j+Q_j)/2\in\Z$. The fractions $R_j/Q_j$ are the convergents of
$\beta=(\ee+1)/2=[1;1,6,10,14,\dots]$, whose initial convergent pairs are
$(1,1)$ and $(2,1)$ and whose remaining partial quotients are $4j+2$ for
$j\ge1$. Applying the Kraaikamp--Liardet criterion \cite{KL} (in the form
recorded on p.~106 of \cite{BDS}) directly to $\beta$ gives
$\|Q_jx\|>\tfrac14 Q_j|Q_j\beta-R_j|$ infinitely often for $x\notin L$. Since
the continued-fraction error formula gives $jQ_j|Q_j\beta-R_j|\to1/4$, this
black-box application yields $\Lam(x)\ge1/16$ outside $L$, which already
recovers the zero-level classification $\{\Lam=0\}=L$, whereas the direct
three-term recurrence analysis in Theorems~\ref{thm:gap}
and~\ref{thm:full-spectrum} yields the sharp initial gap $1/4$ and the complete
quantization $\Lam(\R)=\frac14\Z_{\ge0}\cup\{\infty\}$.
\end{itemize}
For a general irrational $\alpha$ whose partial quotients tend to infinity, we
First treat polynomial-growth denominator recurrences
$Q_{j+1}=(cj^r+d)Q_j+Q_{j-1}$ in
Theorems~\ref{thm:general-rec}--\ref{thm:poly-phase-transition} and then, in
Section~\ref{sec:growing-rec} (Theorems~\ref{thm:growing-rec-spectrum}
and~\ref{thm:growing-rec-dim-phase}, and
Corollary~\ref{cor:growing-rec-three-regimes}), establish the complete unified
theory for arbitrary growing-quotient recurrences $Q_{j+1}=A_jQ_j+Q_{j-1}$ with
$A_j\in\Z_{\ge2}$ and $A_j\to\infty$. In particular, Theorem~\ref{thm:growing-rec-spectrum}(i)
proves the exact asymptotic identity
$\limsup_{j\to\infty}\|Q_jx\|/(Q_j|Q_j\theta-P_j|)=\limsup_{j\to\infty}A_j\|Q_jx\|$
matching the Kraaikamp--Liardet error normalization across all growing-quotient
recurrences, while Theorem~\ref{thm:growing-rec-dim-phase} reveals how the upper
logarithmic growth exponent $\gamma_{\sup}=\limsup_{j\to\infty}\frac{\log A_j}{\log Q_j}\in[0,\infty]$
governs both the subcritical Hausdorff-dimension distortion
$(1-\tau)/(1+\tau\gamma_{\sup})$ ($0\le\tau<1$) and the supercritical secondary
collapse from $L_G$ to $\frac{1}{g}\Z$ across $\tau_{\mathrm{crit}}=1+1/\gamma_{\sup}$.
For the modular continued-fraction alphabets $F_{m,r}=\{[0;a_1,a_2,\dots]:a_n\equiv r\pmod m\}$,
infinite-alphabet thermodynamic-formalism dimension theory goes back to Hensley,
Jenkinson--Pollicott, and Chousionis--Leykekhman--Urba\'nski \cite{CLU}; here
Theorem~\ref{thm:residues-asymptotic} combines an exact closed-form resolvent
identity for the truncated metallic-continuant transfer operator with our
unconditional transfer-operator sandwich to prove $\dimH F_{m,m}<\dimH F_{m,1}=\dimH E_m$
for every $m\ge2$ and to extract the explicit two-term asymptotic expansion
$\dimH F_{m,r}=\frac12+\frac{C_r}{2m}-\frac{C_r^2\log m}{2m^2}+O(m^{-2})$ governed
by the reciprocal metallic-continuant series $C_r=\sum_{k=1}^\infty 1/q_k(r)$ and,
at $r=1$, the reciprocal Fibonacci constant $\psi_{\mathrm{Fib}}=\sum_{k=1}^\infty 1/F_k$
(a targeted negative search of arXiv and prior continued-fraction dimension
literature found no prior appearance of $\psi_{\mathrm{Fib}}$ or $C_r$ as the
asymptotic dimension constant of modular continued-fraction digit sets).
Accordingly, we
regard Theorems~\ref{thm:class} and~\ref{thm:gap} as a variant of known results.
The main new contributions are the complete spectrum discretization and
three-way positive-level splitting (Theorems~\ref{thm:full-spectrum},
\ref{thm:trichotomy}, \ref{thm:general-rec}--\ref{thm:poly-phase-transition},
and~\ref{thm:growing-rec-spectrum}--\ref{thm:growing-rec-dim-phase}), the critical gauge theorems
(Theorem~\ref{thm:gauge} and Corollary~\ref{cor:exact-gauge}), the machine-checked
dimension-zero and subcritical dimension theorems (Theorems~\ref{thm:finite},
\ref{thm:upper}, and~\ref{conj:weight}), and the continued-fraction window and
residue theorems (Theorems~\ref{thm:windows}--\ref{thm:residues-asymptotic}).

Related metric questions for Cantor series expansions, where the approximation
is measured by the tail $x-\omega_n(x)$ of the partial sums, were studied by
Cheng and Zhang~\cite{CZ}, who determine Hausdorff dimensions of the
corresponding limsup sets; spectra and exact gauges are not considered there.
For the selected sequence $(Q_j)$, Corollary~8 of Bugeaud, Liao and Rams
\cite{BLR} applies directly since $Q_{j+1}/Q_j\to\infty$. Theorems~\ref{thm:poly-phase-transition} and
\ref{thm:growing-rec-dim-phase} are compared with \cite{BLR} in
Remarks~\ref{rem:poly-BLR} and~\ref{rem:growing-BLR}. It gives dimension
one for every constant target $c\in[0,1/2]$ at weight $\alpha=0$ in
Theorem~\ref{conj:weight} (in particular for $\{x:\|Q_jx\|\to0\}$), by taking
the centres equal to $c$ and using $|\|y\|-\|c\||\le\|y-c\|$; integer
translation invariance transfers full dimension between $\R$ and $[0,1]$. For
nonzero weights $0<\alpha<1$, the logarithmic growth condition (1.6) of their
Theorem~7 fails for the linear maps $x\mapsto Q_jx$ because
$\log Q_{j+1}/\log Q_j\to1$, and classical mass-transference methods for
$\limsup$ shrinking-target sets \cite{BV} do not apply to the eventual
(liminf-type) bounds defining $E_c$ and the prescribed-limit sets. We establish
the exact dimension $1-\alpha$ here by an elementary covering upper bound and a
shifted-grid Moran construction.

Regarding Hausdorff gauges at critical levels, Cabrelli, Mendivil, Molter and
Shonkwiler \cite[\S5, Theorem~3]{CMMS} construct a gauge giving positive
\emph{finite} measure to a Cantor set from its gap sequence, whereas Moreira,
Erazo and Angelini \cite{MEA} show that the classical Lagrange level at value
$3$ has either zero or non-$\sigma$-finite measure for every dimension function.
Here the recurrence specifies the gauge $h_k(1/Q_n)=(2k+1)^{-n}$, and the entire
critical level $\{\Lam=k/4\}$ on $\R/\Z$ has \emph{infinite, $\sigma$-finite}
$\mathcal H^{h_k}$ measure, proved by combining a countable decomposition by
initial recurrence data with translation invariance under the dense subgroup
$L$.

Relative to earlier literature, our setting has four distinguishing features:
(i) a \emph{subsequence} of convergent denominators of $\ee$ (every third),
(ii) the critical \emph{weight} $j$, (iii) the lattice
$L=\Z+\Z\frac{\ee+1}{2}$ (note $(\ee+1)/2$, not $\ee$), and (iv) the
\emph{sharp} quantum $1/4$. The weight $j$ is \emph{critical} for the covering
argument: it makes the per-parent branching number uniformly bounded, yielding
dimension $0$ for \emph{all} finite level sets (Theorem~\ref{thm:finite}),
whereas without the weight the branching grows linearly and the unweighted
zero-set $\{x:\|Q_jx\|\to0\}$ has full dimension $1$, while the weighted
zero-set $\{x:j\|Q_jx\|\to0\}$ equals the countable group $L$.

